% ============================================================
%  sections/introduction.tex  —  Introduction (merged & cleaned)
%  Sources: SOC draft + sec/1_intro.tex (previous writeup)
% ============================================================
\section{Introduction}
\label{sec:introduction}
%Autonomous driving is something that people are actively exploring, and there existws a lot of simulators for autonomous vehicles. Why researchers need simulators for AV? Because we surely need a system that is not real world, to try the trained jdriving algorithms, as the real world deployment is costly, it's impossible to try things in scale in the real world. Local transportation departments are going to make a lot of administrative decisions every day, especially for urban areas. I should take some data about how much traffic decisions are made in a day, and there must be researches about the disruptive effect to traffic from those decisions. Traffic engineers spent decades to make traffic simulators that work on both microscopic level and macroscopic level. ON one hand, we have SUMO, which is an opensource traffic simulator, that is widely used in digital twin technology. SUMO is a quite macroscope simjulator, which simulates how the traffic system will respond to a certain change to the system, such as road closure, lane closure. However, there are plenty of things that are simplified in SUMO, such as vehicle dynamics. Differetn weathers will have effects on how cars dynamics change, and how the drivers behaviors change. Those changes cannot be captured by SUMO's modeling of driver model, since SUMO does not take environmental conditions into account. There exists simulators that do count those factors, such as Carla \cite{dosovitskiy_carla_2017}who's built on Unreal engine, the backend is PhysX physics engine. Also metadrive \cite{li_metadrive_2022}, who built their system on bullet physics engine. However, one shared downside of them is that both simulators, although enabled by physics engine, they don't support scaling of scenes, for example, if i would like to evaluate the potentail outcome of a certain configuration of a workzone, I need to sequentially test all of them. This is requiring a lot of time. There's a gap in current AV siimulators, that that can both be scaled to multiple scene, multiple agents at one time, and also have phyiscs engine enabled.

% ------------------------------------------------------------------
%  P1 — Broad motivation: why simulation matters for AV
% ------------------------------------------------------------------

Simulators are indispensable to the development of autonomous driving.
Learning-based driving agents require platforms where their behavior can be evaluated in a controlled, reproducible, and safe manner—especially for rare or dangerous scenarios that cannot be ethically staged on public roads~\cite{grigorescu_av_survey_2020}.
There is a well-known gap between simulators and the real world, and researchers have been actively pursuing better simulation to narrow it.
A fundamental tension exists in simulator design: higher reconstruction fidelity—whether in physics, visuals, or both—demands greater computational cost, while systems that prioritize scalability, such as Waymax~\cite{gulino_waymax_2023} and GPUDrive~\cite{kazemkhani_gpudrive_2025}, simplify the underlying dynamics to bicycle models or other parametrized representations.
Recent GPU-parallel computing technologies have begun to shift this trade-off by moving heavy computation onto graphics processing units, opening the door to large-scale benchmarking for autonomous driving.
As Kalra and Paddock demonstrated, hundreds of millions to billions of miles of driving would be needed for a single autonomous vehicle to statistically demonstrate its reliability in terms of fatalities and injuries~\cite{kalra_driving_safety_2016}—a scale that is infeasible through real-world testing alone and underscores the necessity of high-throughput simulation.

% ------------------------------------------------------------------
%  P2 — Weather and safety motivation (from previous draft)
% ------------------------------------------------------------------

The safety relevance of simulation is further amplified by real-world driving conditions.
According to the Federal Highway Administration (FHWA)~\cite{noauthor_how_nodate}, approximately 6{,}000{,}000 vehicle crashes occur annually in the United States, of which roughly 12\% are weather-related.
Among those, 94\% involve changes in pavement friction: 73\% occur during rain or mist and 21\% during frozen precipitation.
Human driving behavior naturally adapts to adverse weather—drivers become more cautious when rainfall is high, compensating for reduced visibility and diminished maneuverability as effective road friction decreases~\cite{chen_influence_2019, ahmed_impacts_2018, peng_examining_2018}.
However, most existing simulation frameworks do not explicitly capture this behavioral adaptation, which may lead to unrealistic driving behavior in evaluation.
Accurate modeling of weather-dependent vehicle dynamics is particularly important for transportation planning, traffic safety analysis, and risk assessment of road infrastructure under variable conditions.

% ------------------------------------------------------------------
%  P3 — Traffic management decision-making motivation
% ------------------------------------------------------------------

Beyond autonomous-vehicle development, simulators also serve traffic administrators who make high-stakes operational decisions daily—road closures, lane reconfigurations, work zone designs—whose effects propagate into system-level economic and safety costs.
According to FHWA work zone operations guidelines, agencies actively monitor delay, queue length, and safety indicators, iteratively revising traffic management plans during execution~\cite{fhwa_workzone_2023}.
With the emergence of high-throughput simulators, administrators are increasingly moving from observing real system responses to evaluating candidate policies in simulation before deployment.
Yet, without accurate modeling of scene physics, ground friction, and plausible actor behavior, simulation-based decisions cannot be trusted, and suboptimal or dangerous policies may be selected.

% ------------------------------------------------------------------
%  P4 — Current simulators and their limitations
% ------------------------------------------------------------------

Traffic engineers have spent decades developing simulators at both microscopic and macroscopic levels.
SUMO~\cite{krajzewicz_sumo_2012}, for instance, is a widely used open-source microscopic traffic simulator employed in digital-twin applications, tracking individual vehicles with car-following and lane-changing models.
However, SUMO provides only limited support for weather effects and does not model tire--road friction physics.
Simulators such as CARLA~\cite{Dosovitskiy17}, built on Unreal Engine~4 with PhysX-based vehicle physics, and MetaDrive~\cite{li2022metadrive}, built on Panda3D with the Bullet engine, incorporate richer vehicle dynamics and environmental modeling.
However, in these systems each vehicle is driven by non-vectorized per-agent Python loops: actuation commands, state queries, and reward computations are executed sequentially for every agent, making the CPU-side control code---not the physics solver---the throughput bottleneck.
As a result, neither simulator supports \emph{GPU-parallel execution of multiple independent scenes}: evaluating $N$ scenarios requires $N$ sequential simulation runs, which is prohibitively slow for large-scale benchmarking or policy search.

On the other end of the spectrum, GPU-accelerated simulators such as Waymax~\cite{gulino_waymax_2023} and GPUDrive~\cite{kazemkhani_gpudrive_2025} achieve impressive throughput by batching hundreds of scenarios on a single GPU.
GPUDrive, built on the Madrona game engine, can simulate over a million steps per second for multi-agent training on the Waymo Open Motion Dataset.
However, these systems achieve scalability by simplifying the underlying dynamics. The vehicles follow bicycle or kinematic models with no rigid-body physics engine governing tire forces, suspension, or contact dynamics.
This simplification is reasonable for many research tasks, but limits applicability to safety-critical analyses where physical fidelity matters.

% ------------------------------------------------------------------
%  P5 — The gap statement
% ------------------------------------------------------------------

There exists a gap in the current landscape of autonomous-driving simulators: no existing system simultaneously provides
(i)~GPU-parallel execution of multiple independent driving scenes,
(ii)~multi-agent support within each scene,
(iii)~\emph{vectorized} articulated vehicles whose actuation, sensing, and reward computation scale with GPU parallelism rather than with per-agent CPU loops, and
(iv)~a full rigid-body physics engine governing vehicle dynamics with weather-aware friction modeling.

% ------------------------------------------------------------------
%  P6 — Our solution
% ------------------------------------------------------------------

To fill this gap, we present \textbf{SceneFactory}, a procedural scene construction and multi-agent training platform built on NVIDIA Isaac Sim and Isaac Lab~\cite{nvidia_isaac_2025}.
Isaac Lab provides a vectorized multi-agent reinforcement learning environment (DirectMARLEnv) that runs entirely on the GPU, with PhysX~5 as the underlying rigid-body physics engine.
A key technical contribution is a \emph{fully vectorized articulated PhysX vehicle}: all joint-state queries, actuation commands, observation assembly, and reward computation are expressed as batched tensor operations over the full agent population, eliminating the per-agent Python loops that throttle existing physics-based simulators.
This vectorized design achieves up to \textbf{127$\times$} higher throughput than a non-vectorized PhysX baseline on the same GPU and physics solver (\S\ref{sec:throughput}).
SceneFactory ingests real-world road topologies from the Waymo Open Motion Dataset~\cite{Kan_2024_icra}, procedurally constructs physically realistic USD road environments, and orchestrates the parallel simulation of hundreds of these articulated vehicles across dozens of distinct road worlds on a single GPU.
Additionally, we integrate a weather-to-friction module based on Zhao et al.'s tire--pavement friction model~\cite{zhao_tire-pavement_2024}, which maps precipitation conditions and road surface type to effective friction coefficients applied directly to PhysX material properties.
This paper presents the first phase of a larger research program; it describes the full pipeline and reports initial training results.

% ------------------------------------------------------------------
%  P7 — Contribution list
% ------------------------------------------------------------------

The contributions of this work are:
\begin{enumerate}[leftmargin=*, itemsep=2pt]
  \item \textbf{SceneFactory}, a procedural scene construction pipeline that converts Waymo Open Motion Dataset scenarios into GPU-parallelized, multi-world Isaac Sim stages with full PhysX rigid-body dynamics.
  \item A \textbf{fully vectorized articulated PhysX vehicle} whose joint-state queries, actuation, observation assembly, and reward computation are expressed entirely as batched GPU tensor operations, achieving up to 127$\times$ higher throughput than a non-vectorized PhysX baseline on the same hardware.
  \item A \textbf{multi-agent RL training environment} supporting $N$ parallel worlds $\times$ $M$ agents per world, all driven by a shared navigation policy trained via PPO.
  \item A \textbf{weather-to-friction module} that maps precipitation and road surface type to PhysX tire friction coefficients, enabling weather-aware driving behavior training.
\end{enumerate}

\chattodo{Insert concrete throughput and success-rate numbers once experiments are finalized (e.g., ``X env-steps/sec with Y agents on a single A6000 pro, Z\% goal-reach rate'').}
